\chapter{Un cuerpo que habla}

\epigraph{El primer lenguaje que aprende un ser humano no tiene palabras. Se habla con la piel.}{}

En el capítulo anterior vimos que nacemos radicalmente inmaduros ---un feto externo, como decía Montagu--- y que esa vulnerabilidad no es un fallo sino el precio de nuestro cerebro extraordinario. Vimos también, casi de pasada, que la evolución no nos dejó solos ante esa fragilidad: preparó mecanismos que empujan a cuidar.

Ahora toca detenernos en esos mecanismos. Porque lo que ocurre en el cuerpo de una madre durante el parto, lo que le pasa al padre cuando sostiene a su hijo por primera vez, lo que sucede entre la piel de un recién nacido y la piel de quien lo abraza, es mucho más que biología. Es un lenguaje. Un lenguaje que nuestro cuerpo habla antes de que nosotros hayamos aprendido una sola palabra.

\section{La inundación}

En las últimas fases del parto, el cuerpo de la madre hace algo extraordinario. El hipotálamo, una estructura pequeña y antigua en la base del cerebro, comienza a liberar oxitocina en cantidades masivas.\footnote{Los niveles de oxitocina en plasma durante el parto pueden ser tres o cuatro veces superiores a los niveles basales, con picos aún mayores en el momento de la expulsión. Cf. Uvnäs-Moberg, K., \textit{The Oxytocin Factor}, Da Capo Press, 2003.} No es una liberación gradual. Es una inundación.

La oxitocina hace que el útero se contraiga ---esa es su función más conocida, la que aprovecha la medicina cuando induce un parto---. Pero hace mucho más. Actúa en el cerebro como un interruptor emocional. Reduce la ansiedad. Genera una sensación profunda de calma y confianza. Y, sobre todo, activa los circuitos cerebrales del vínculo: esos que hacen que una madre mire a la criatura que acaba de salir de su cuerpo y sienta algo que ninguna palabra describe del todo. Algo que no es exactamente decisión, ni exactamente instinto, sino una mezcla poderosa de ambos.

El torrente no se detiene con el parto. La lactancia lo mantiene. Cada vez que el bebé succiona, se libera más oxitocina.\footnote{La succión del pezón estimula las terminaciones nerviosas que envían señales al hipotálamo, desencadenando la liberación pulsátil de oxitocina. Este mecanismo tiene una doble función: provoca la eyección de leche (el ``reflejo de bajada'') y simultáneamente refuerza el vínculo emocional madre-hijo. Cf. Rilling y Young, ``The biology of mammalian parenting and its effect on offspring social development'', \textit{Science}, 2014.} No solo para que la leche fluya ---que también---, sino para que el vínculo se refuerce. La naturaleza, si se me permite la expresión, no se fía de una sola dosis. Insiste. Cada toma es un recordatorio biológico: \textit{esta criatura es tuya, cuídala}.

Y hay algo más. Durante el parto y la lactancia también se liberan endorfinas ---los opiáceos naturales del cuerpo--- que generan una sensación de bienestar profundo.\footnote{Cf. Matthiesen et al., ``Postpartum maternal oxytocin release by newborns: effects of infant hand massage and sucking'', \textit{Birth}, 2001.} Es como si el cuerpo recompensara a la madre por el esfuerzo del parto y por el acto de amamantar. Como si le dijera: \textit{esto que estás haciendo es lo correcto, sigue así}.

No es magia. Es química. Pero es una química que apunta con una precisión asombrosa hacia algo que todos reconocemos como profundamente humano: el amor de una madre por su hijo.

\section{La piel que siente}

Hay un momento que los obstetras conocen bien. Inmediatamente después del parto, si se coloca al recién nacido desnudo sobre el pecho desnudo de su madre ---lo que se llama contacto piel con piel---, ocurre algo notable.

El bebé, que hace un instante estaba llorando, se calma. Su ritmo cardíaco se estabiliza. Su temperatura corporal se regula ---la piel de la madre ajusta su temperatura para calentar o enfriar al bebé según lo necesite---.\footnote{Este fenómeno se conoce como ``sincronía térmica''. El pecho de la madre puede variar hasta dos grados centígrados para ajustarse a las necesidades térmicas del recién nacido. Cf. Ludington-Hoe, S., ``Skin-to-skin contact: A calming and reassuring experience for preterm infants'', \textit{Neonatal Network}, 2004.} Sus niveles de cortisol ---la hormona del estrés--- bajan drásticamente. Si se le deja ahí, sin prisa, el bebé es capaz de reptar lentamente hacia el pecho y encontrar el pezón por sí solo, guiado por el olfato.\footnote{Este comportamiento, documentado en vídeo por investigadores suecos en los años noventa, se conoce como \textit{breast crawl}. El recién nacido, colocado sobre el abdomen de la madre, es capaz de desplazarse hacia el pecho y comenzar a mamar sin ninguna ayuda externa en los primeros 60-90 minutos de vida. Cf. Widström et al., ``Newborn behaviour to locate the breast when skin-to-skin'', \textit{Acta Paediatrica}, 2011.}

No es solo el bebé el que responde. La madre también. El contacto piel con piel dispara en ella una nueva oleada de oxitocina. Activa circuitos cerebrales que los neurocientíficos asocian con la recompensa y la motivación.\footnote{Estudios con resonancia magnética funcional muestran que el contacto con el recién nacido activa en la madre el área tegmental ventral y el núcleo accumbens ---los mismos circuitos que se activan con otras experiencias de recompensa intensa---. Cf. Strathearn et al., ``Adult attachment predicts maternal brain and oxytocin response to infant cues'', \textit{Neuropsychopharmacology}, 2009.} En términos llanos: tocar a su bebé le produce placer. Un placer que no es frívolo sino profundamente funcional, porque la empuja a seguir en contacto, a seguir cuidando.

La ciencia ha medido estos efectos con rigor, y los resultados son consistentes: los bebés que reciben contacto piel con piel inmediato lloran menos, duermen mejor, mantienen niveles de glucosa más estables y tienen más éxito con la lactancia. Los prematuros que reciben el llamado ``método canguro'' ---contacto piel con piel prolongado--- ganan peso más rápido, regulan mejor la temperatura y tienen menos infecciones.\footnote{El método canguro fue desarrollado en Colombia en 1978 por los doctores Rey y Martínez como respuesta a la escasez de incubadoras. Los resultados fueron tan buenos que hoy la OMS lo recomienda como estándar de cuidado para prematuros en todo el mundo. Una revisión Cochrane de 2016 concluyó que reduce la mortalidad neonatal en un 40\%. Cf. Conde-Agudelo y Díaz-Rossello, ``Kangaroo mother care to reduce morbidity and mortality in low birthweight infants'', \textit{Cochrane Database of Systematic Reviews}, 2016.}

La piel, ese órgano enorme que nos envuelve, resulta ser mucho más que una barrera protectora. Es un órgano de comunicación. El primer canal por el que un ser humano recibe un mensaje, antes de entender una sola palabra, antes de reconocer una cara, es la piel. Y el mensaje que recibe es: \textit{estás a salvo, hay alguien aquí}.

\section{El cerebro del padre}

Pero este capítulo quedaría incompleto ---y sería injusto--- si se limitara a la madre. Porque uno de los hallazgos más reveladores de la neurociencia reciente es que el cerebro del padre también cambia.

Ya lo mencionamos brevemente en el capítulo anterior: cuando un hombre se convierte en padre, sus niveles de testosterona bajan y los de oxitocina suben. Pero el cambio va más allá de las hormonas. Estudios con neuroimagen muestran que la paternidad reorganiza literalmente el cerebro masculino.\footnote{Cf. Abraham et al., ``Father's brain is sensitive to childcare experiences'', \textit{PNAS}, 2014. Este estudio comparó la actividad cerebral de madres, padres heterosexuales y padres homosexuales primarios, y encontró que la implicación activa en el cuidado ---más que el sexo biológico--- era lo que determinaba la activación de los circuitos parentales.}

En los padres que se implican en el cuidado, se activan regiones asociadas con la empatía, la lectura emocional y la vigilancia. La amígdala ---que procesa las amenazas--- se sensibiliza al llanto del bebé. El córtex prefrontal ---que planifica y toma decisiones--- se conecta más intensamente con las áreas emocionales. Es como si el cerebro del hombre se recableara para incorporar una nueva prioridad: \textit{proteger a esta criatura}.

Y hay un dato que me parece especialmente elocuente. La bajada de testosterona no es un accidente ni un efecto secundario indeseable. Es proporcional a la implicación: cuanto más tiempo pasa un padre con su bebé, más baja la testosterona.\footnote{Cf. Gettler et al., ``Longitudinal evidence that fatherhood decreases testosterone in human males'', \textit{PNAS}, 2011. Los padres que dormían en la misma habitación que sus hijos mostraban los niveles más bajos de testosterona.} No es que ser padre ``debilite'' al hombre, como cierta cultura querría hacernos creer. Es que lo reconfigura. La testosterona alta es útil para competir, para buscar pareja, para enfrentarse a rivales. Pero cuando ya tienes pareja y tienes una cría vulnerable, lo que necesitas no es competición: es presencia, atención, cuidado. Y el cuerpo lo sabe.

Hay aquí algo que debería darnos que pensar. El cuerpo del hombre no está diseñado solo para engendrar. Está diseñado también para quedarse. Para cuidar. Para ser padre, no solo progenitor. La biología no nos convierte en padres automáticamente ---eso sería determinismo, y los seres humanos no somos máquinas---. Pero nos pone el viento a favor. Nos da un empujón en la dirección correcta, igual que se lo da a la madre.

\section{Un cuerpo configurado para la relación}

Si damos un paso atrás y miramos el panorama completo, lo que encontramos es impresionante.

Tenemos un parto que libera hormonas del vínculo. Una lactancia que las mantiene. Una piel que, al tocar otra piel, dispara cascadas de bienestar y confianza. Un cerebro materno que se reorganiza para la vigilancia y el cuidado. Un cerebro paterno que hace lo mismo cuando el padre se implica. Un bebé cuyo llanto activa circuitos de alarma en los adultos de su alrededor. Un rostro infantil ---ojos grandes, frente ancha, mejillas redondas--- que dispara ternura en prácticamente todos los seres humanos.\footnote{El etólogo Konrad Lorenz describió en 1943 el \textit{Kindchenschema} (``esquema infantil''): un conjunto de rasgos faciales que los humanos percibimos como ``tiernos'' y que activan conductas de cuidado. Estudios posteriores han confirmado que estos rasgos provocan respuestas cerebrales de recompensa incluso en personas que no son padres. Cf. Lorenz, K., ``Die angeborenen Formen möglicher Erfahrung'', \textit{Zeitschrift für Tierpsychologie}, 1943.}

Todo esto no es casualidad. Es un sistema. Un sistema en el que cada pieza encaja con las demás para producir un resultado: que los adultos cuiden de una cría que no puede cuidarse sola.

Podemos explicarlo en términos evolutivos, y la explicación es sólida: los grupos donde estos mecanismos funcionaban mejor criaban más hijos que llegaban a la edad adulta, y esos hijos transmitían los mismos mecanismos. Selección natural. Pero la explicación evolutiva, siendo verdadera, no agota el significado de lo que vemos. Porque lo que vemos es un cuerpo que no está diseñado para la soledad. Un cuerpo que habla, con cada una de sus fibras, un lenguaje de relación.

Y ese lenguaje no se limita a la relación madre-hijo o padre-hijo. Como vamos a ver ahora, marca para siempre la forma en que nos relacionamos con todos los demás.

\section{Bowlby y el descubrimiento del apego}

En 1950, la Organización Mundial de la Salud encargó a un psiquiatra británico llamado John Bowlby un informe sobre la salud mental de los niños sin hogar en la Europa de posguerra.\footnote{Bowlby, J., \textit{Maternal Care and Mental Health}, OMS, 1951. Este informe fue el germen de lo que después se convertiría en la teoría del apego, desarrollada en su trilogía \textit{Attachment and Loss} (1969, 1973, 1980).} Lo que Bowlby encontró cambió para siempre nuestra forma de entender el desarrollo humano.

Miles de niños habían sido separados de sus padres durante la guerra ---evacuados, huérfanos, internados en instituciones---. Muchos de ellos habían recibido alimento suficiente, abrigo, atención médica. Pero algo les faltaba. No crecían bien. Algunos dejaban de hablar, se balanceaban compulsivamente, se negaban a comer. Otros se mostraban incapaces de establecer relaciones con los cuidadores. Las tasas de mortalidad en los orfanatos eran alarmantemente altas, incluso cuando las condiciones materiales eran aceptables.\footnote{El pediatra René Spitz documentó este fenómeno en su investigación sobre lo que llamó ``hospitalismo'': niños institucionalizados que, pese a recibir alimentación y cuidado médico adecuados, desarrollaban un deterioro físico y psicológico severo por falta de contacto humano. Algunos morían. Cf. Spitz, R., ``Hospitalism: An inquiry into the genesis of psychiatric conditions in early childhood'', \textit{The Psychoanalytic Study of the Child}, 1945.}

Bowlby comprendió algo que hoy nos parece evidente pero que entonces era revolucionario: \textbf{un niño no necesita solo alimento y protección. Necesita una relación}. Necesita una figura estable, predecible, que responda a sus llamadas. Necesita lo que Bowlby llamó una \textit{base segura}: alguien a quien agarrarse cuando tiene miedo, desde cuya cercanía puede explorar el mundo.

Y esto no era un lujo emocional. Era una necesidad biológica. Bowlby lo planteó en términos evolutivos: en la sabana donde evolucionamos, un bebé separado de su cuidador era un bebé muerto. El llanto, la búsqueda de proximidad, la angustia ante la separación, no eran caprichos ni ``mañas''. Eran un sistema de supervivencia, tallado por millones de años de selección natural.

Ese sistema es lo que hoy llamamos \textbf{apego}.

\section{Lo que aprendemos sin saber que lo aprendemos}

Unos años después de Bowlby, una psicóloga canadiense llamada Mary Ainsworth ideó un experimento elegante para observar el apego en acción.\footnote{El procedimiento, conocido como ``Situación Extraña'' (\textit{Strange Situation}), fue diseñado por Ainsworth a finales de los años sesenta. Se basa en observar la reacción de un niño de entre 12 y 18 meses ante una secuencia de separaciones y reencuentros con su madre en presencia de un desconocido. Cf. Ainsworth et al., \textit{Patterns of Attachment}, Lawrence Erlbaum, 1978.}

La situación era sencilla: una madre entraba con su hijo pequeño en una habitación con juguetes. Un desconocido entraba. La madre salía. El desconocido intentaba consolar al niño. La madre volvía. Los investigadores observaban las reacciones del niño en cada fase.

Lo que Ainsworth descubrió fue que no todos los niños reaccionaban igual. Y que las diferencias no eran aleatorias: tenían que ver con la calidad de la relación con la madre durante los meses anteriores. Ainsworth identificó tres patrones principales:

\subsection{Apego seguro}

La mayoría de los niños ---alrededor del 60-65\%--- mostraban lo que Ainsworth llamó \textit{apego seguro}. Estos niños usaban a la madre como base para explorar: jugaban con confianza mientras ella estaba presente, se angustiaban al verla partir, y al volver a verla, iban hacia ella, se consolaban rápidamente y volvían a jugar.

Estos niños tenían madres que, en general, respondían a sus necesidades de forma consistente: cuando el bebé lloraba, acudían; cuando buscaba contacto, lo daban; cuando necesitaba espacio, lo permitían. No madres perfectas ---eso no existe---. Madres suficientemente atentas.\footnote{El concepto de ``madre suficientemente buena'' (\textit{good enough mother}) fue acuñado por el pediatra y psicoanalista Donald Winnicott, y refleja una idea liberadora: no se necesita la perfección para criar bien. Se necesita presencia, atención y una respuesta razonablemente consistente. Cf. Winnicott, D. W., \textit{Playing and Reality}, Tavistock, 1971.}

\subsection{Apego inseguro}

Otros niños mostraban patrones distintos. Algunos, llamados \textit{evitativos}, parecían indiferentes a la partida de la madre y la ignoraban al volver. No es que no les importara ---sus niveles de cortisol estaban tan altos como los de los demás---, sino que habían aprendido que expresar la necesidad no servía de nada. Sus madres tendían a ser menos receptivas al contacto físico, más distantes emocionalmente.

Otros, llamados \textit{ambivalentes} o \textit{ansiosos}, se mostraban angustiadísimos cuando la madre se iba, pero al volver no se consolaban: se aferraban a ella con furia, llorando, resistiéndose a ser calmados. Habían aprendido que la respuesta de la madre era impredecible ---a veces atenta, a veces ausente--- y habían desarrollado una estrategia de hiperactivación: gritar más fuerte, llorar más tiempo, para asegurarse de que alguien respondiera.\footnote{Investigaciones posteriores añadieron un cuarto patrón, el apego \textit{desorganizado}, observado en niños cuyo cuidador era al mismo tiempo fuente de consuelo y fuente de miedo (como en casos de maltrato). Cf. Main y Solomon, ``Discovery of a new, insecure-disorganized/disoriented attachment pattern'', en Brazelton y Yogman (eds.), \textit{Affective Development in Infancy}, Ablex, 1986.}

\section{Del bebé al adulto}

Aquí viene lo que de verdad importa para quienes estáis leyendo este libro pensando en vuestro matrimonio.

Los estilos de apego no se quedan en la infancia. Nos acompañan. Investigaciones de las últimas décadas han mostrado que los patrones que aprendemos en los primeros años ---esa forma de estar con el otro, esa expectativa sobre si los demás van a responder cuando los necesitemos--- tienden a reproducirse en las relaciones adultas, incluida la relación de pareja.\footnote{Cf. Hazan y Shaver, ``Romantic love conceptualized as an attachment process'', \textit{Journal of Personality and Social Psychology}, 1987. Este artículo pionero trasladó la teoría del apego al ámbito de las relaciones adultas y abrió una de las líneas de investigación más fructíferas de la psicología de la pareja.}

Una persona con apego seguro tiende a sentirse cómoda con la intimidad. Puede pedir ayuda sin sentirse débil. Puede dar espacio al otro sin sentirse abandonada. Confía, en general, en que el otro va a estar ahí.

Una persona con apego evitativo tiende a la independencia excesiva. Le cuesta pedir, le cuesta mostrar vulnerabilidad, se retira emocionalmente cuando las cosas se ponen intensas. No es que no necesite al otro; es que aprendió, muy pronto, que mostrar la necesidad era arriesgado.

Una persona con apego ansioso tiende a la hipervigilancia emocional. Necesita confirmación constante de que el otro la quiere. Un mensaje que tarda en llegar se convierte en fuente de angustia. La distancia se interpreta como rechazo. No es que sea ``dependiente'' por elección; es que su sistema de apego aprendió que la disponibilidad del otro era incierta, y se puso en alerta permanente.

Pero ---y esto es fundamental--- \textbf{el apego no es una condena}. Los estilos de apego son tendencias, no destinos. Se pueden conocer, comprender y, con esfuerzo y con la relación adecuada, modificar. De hecho, una relación de pareja segura es uno de los contextos más poderosos para sanar heridas de apego.\footnote{El concepto de ``apego ganado'' (\textit{earned secure attachment}) describe precisamente a personas que, pese a haber tenido experiencias de apego inseguro en la infancia, desarrollan un estilo seguro en la edad adulta ---a menudo a través de una relación de pareja estable, terapia, u otras relaciones significativas---. Cf. Roisman et al., ``Earned-security in retrospect'', \textit{Development and Psychopathology}, 2002.} Cuando alguien te responde de forma constante, cuando está ahí día tras día, tu sistema de apego ---ese sistema antiguo, tallado para la supervivencia--- empieza a relajarse. Empieza a confiar. Empieza a creer que puede soltar la guardia.

\section{Lo que el cuerpo nos dice}

Si recorremos todo lo que hemos visto en este capítulo ---la oxitocina del parto, la piel que comunica, el cerebro del padre que se transforma, el sistema de apego que nos configura--- lo que emerge es una imagen coherente y asombrosa.

Nuestro cuerpo no es un accidente. No es un vehículo neutro que nos transporta por la vida. Es un cuerpo que habla. Que dice, con cada mecanismo, con cada hormona, con cada circuito neuronal: \textit{estás hecho para la relación. Estás hecho para el vínculo. Estás hecho para cuidar y para ser cuidado}.

La biología no es todo, por supuesto. Si lo fuera, bastaría con la química para explicar el amor, y sabemos que no basta. Hay algo en el amor que desborda la biología, algo que le da sentido, algo que la eleva. Algo que hace que una madre que adopta a un niño ---sin el torrente hormonal del parto--- pueda amarlo con la misma intensidad. Algo que hace que un padre al que la biología no le ha puesto las cosas fáciles elija quedarse cada mañana.

Pero la biología es el suelo sobre el que se construye todo lo demás. Y es un suelo que tiene una orientación clara. Como si alguien ---o algo--- hubiera querido que fuera así.

A eso ---a quién o qué hay detrás de ese diseño--- llegaremos más adelante. Por ahora, basta con escuchar lo que el cuerpo dice. Y lo que dice es bastante claro.

\paracaminar

\begin{enumerate}
  \item Conocer los estilos de apego puede ayudar a entenderse mejor como pareja. ¿Os reconocéis en alguno de ellos? ¿Cómo reaccionáis cada uno cuando sentís que el otro se distancia o cuando necesitáis cercanía?

  \item Hablad juntos sobre vuestra experiencia con el contacto físico ---no solo el sexual, sino el cotidiano: la mano en el hombro, el abrazo al llegar, dormir cerca---. ¿Qué papel tiene en vuestra relación? ¿Os resulta fácil o difícil?

  \item Si el cuerpo está ``configurado para la relación'', ¿qué consecuencias tiene eso para la forma en que queréis construir vuestra familia?
\end{enumerate}

\vinetacapitulo{ch02}
